\begin{result}{F.3.34}{An explicit six-vector reconstruction of the Johnson detector}{Proved using the published multiplicity bound and an exact integer certificate}
Let $J_{g,7}$ be the actual rational degree-seven Johnson image
for a genus-$g$ surface with one boundary component, with $g\geq3$.
The five-term tensor map $R$ displayed in (SL.3) satisfies
\[
                              R(J_{g,7})=0.                \tag{RV.1}
\]
The same assertion holds after extension to every $\Q_\ell$.
The proof below reconstructs the specific $R$-calculation omitted
from the sketch in MSS2026 Proposition~4.1. Its dimension input
is the earlier published stable multiplicity bound in the proof
of MSS2015 Theorem~7.5(iv).
\end{result}
\begin{proof}
We exhibit six actual Johnson brackets and use their projected
components to determine the entire part on which $R$ can be nonzero.

\paragraph{The exact older dimension input.}
The proof of MSS2015 Theorem~7.5(iv), printed pp.~29--30,
states that the standard-representation isotypical component
of $J_{G,7}$ has multiplicity six in the stable range
\cite{AuditMoritaSakasaiSuzuki2015v2}. We only need its upper
bound. The source's dimension argument gives multiplicities
$10$ and $3$ for the standard representation in
$\mathfrak h_{G,1}(7)$ and $\mathfrak j_{G,1}(7)$,
respectively, and ranks $3$ and $2$ for the corresponding
ES images. Thus the ES kernel has multiplicity seven and
its intersection with $\mathfrak j_{G,1}(7)$ has multiplicity
one. The source's Proposition~3.5(ii), printed pp.~16--17,
gives $J_{G,7}\cap\mathfrak j_{G,1}(7)=0$; together with
ES vanishing on $J_{G,7}$, this bounds its multiplicity by six.

These are explicitly retained published dimension inputs.
The present certificate does not claim to recompute all
the character and ES-rank calculations in that older proof.
We first choose a genus $G\geq5$ in its stable range, and
then pass to arbitrary genus by the actual extension map
below. No sharper numerical cutoff for the older
degree-seven calculation is needed.

\paragraph{Six actual vectors and six equivariant maps.}
Write $e_0=x_1,e_1=y_1,\ldots,e_8=x_5,e_9=y_5$, with
$\langle x_i,y_i\rangle=1$. For an ordered triple
$t=(i,j,k)$ let
\[
 \delta_t(z)=\langle z,e_i\rangle[e_j,e_k]
            +\langle z,e_j\rangle[e_k,e_i]
            +\langle z,e_k\rangle[e_i,e_j].
\]
Every $\delta_t$ is an actual degree-one Johnson vector.
For each row of the following table form the left-nested
sevenfold bracket
\[
 Z_r=[[[[[[\delta_{t_1},\delta_{t_2}],\delta_{t_3}],
               \delta_{t_4}],\delta_{t_5}],\delta_{t_6}],
               \delta_{t_7}].
\]
The actual Johnson image is a Lie subalgebra, so all six
$Z_r$ lie in $J_{G,7}$.
\[
\begin{array}{c|ccccccc}
r&t_1&t_2&t_3&t_4&t_5&t_6&t_7\\\hline
1&(3,7,5)&(1,2,0)&(9,8,6)&(6,8,5)&(2,4,7)&(9,0,3)&(4,6,7)\\
2&(9,5,0)&(6,4,7)&(6,2,0)&(6,5,8)&(7,3,8)&(1,4,3)&(9,2,7)\\
3&(1,3,4)&(4,2,0)&(5,7,0)&(9,2,7)&(5,6,8)&(5,3,0)&(1,4,6)\\
4&(0,1,4)&(3,5,0)&(9,0,5)&(2,3,9)&(4,6,8)&(8,1,0)&(2,7,1)\\
5&(3,7,0)&(4,9,6)&(3,5,2)&(7,3,8)&(6,7,2)&(0,8,9)&(1,2,6)\\
6&(5,4,0)&(0,8,1)&(5,2,3)&(7,4,3)&(0,4,1)&(6,5,2)&(6,7,9).
\end{array}                                                \tag{RV.2}
\]
Each has torus weight $x_1$. For a matching $C$ of
positions $1,\ldots,8$, define the literal equivariant map
\[
 K_C(u_1\otimes\cdots\otimes u_9)
      =\left(\prod_{(i,j)\in C}\langle u_i,u_j\rangle\right)u_9.
                                                               \tag{RV.3}
\]
Use the following six matchings, in order:
\[
\begin{array}{lll}
(12)(34)(57)(68),&(12)(34)(58)(67),&(12)(35)(47)(68),\\
(12)(36)(47)(58),&(12)(36)(48)(57),&(12)(37)(46)(58).
\end{array}
\]
For every one of these maps the value at $Z_r$ is a
multiple of $x_1$. Their scalar coefficients, with
rows indexed by $r$, are exactly
\[
 A=\begin{pmatrix}
18&12&-2&15&-7&-10\\
0&3&0&-4&4&0\\
15&-24&3&2&-5&6\\
-17&14&-10&11&-11&-4\\
-3&-10&2&-15&3&16\\
2&-16&-1&2&-3&0
\end{pmatrix},\qquad \det A=-203784\ne0.                 \tag{RV.4}
\]
Direct evaluation of the five terms in (SL.3) gives
\[
                         R(Z_r)=0\qquad(1\leq r\leq6).
                                                               \tag{RV.5}
\]

The self-contained program
\texttt{check\_R\_J7\_six\_generators.py} constructs the
six full derivations by literal noncommutative substitution.
It verifies degree seven, the weight $x_1$ on every
nonzero tensor term, and specialness $Z_r(\omega)=0$.
It computes all $630$ coefficients of the $105$ possible
last-slot contractions, verifies the displayed minor over
the integers, and evaluates all six $R$-values exactly.
All coefficients use arbitrary-precision integers; the
determinant is exact. The program reads no stored matrix
or witness file. The source's five-term formula is
attributed to MSS2026\cite[Proposition~4.1, p.~11]{AuditMoritaSakasaiSuzuki2026v1};
its vanishing is established here by (RV.4)--(RV.5)
and the argument that follows.

\paragraph{The actual isotypical projection.}
Let $e_H$ denote the symplectic-equivariant projector
onto the isotypical summand of the standard representation
$H=H_1(\Sigma_{G,1},\Q)$ in $\mathfrak h_{G,1}(7)$.
It preserves the actual $J_{G,7}$. Every map in (RV.3),
and $R$ itself, is equivariant with target $H$ and
therefore kills all the other irreducible summands.
Consequently
\[
 K_C(Z_r)=K_C(e_HZ_r),\qquad R(Z_r)=R(e_HZ_r).
\]
The vectors $Z_r$ need not lie entirely in the $H$-isotypical
summand. Equation (RV.4) proves that their six projected
components $e_HZ_r$ are linearly independent. They have
weight $x_1$, and the published multiplicity bound says
that this weight space in $(J_{G,7})_H$ has dimension
at most six. Thus they form a basis of that weight space.
The symplectic orbits of this basis span the whole
standard-representation isotypical summand.

Equation (RV.5) and equivariance show that $R$ kills
that entire summand. It kills every other irreducible
summand because its target is $H$. Hence $R(J_{G,7})=0$.
Only the six displayed equivariant maps and their
nonzero minor are needed; no assertion about completeness
of the $105$ contraction maps enters this argument.
Likewise, only the six explicit brackets use $J_1$;
the conclusion does not require a new proof that every
element of $J_7$ is generated in degree one.

\paragraph{Every genus and every rational prime.}
The six vectors have support on five handles. Under
the actual symplectic inclusion into a larger genus,
their tensor expressions, the six-map matrix, and
all their $R$-values are unchanged. The preceding
argument therefore applies in every sufficiently
large stable genus.

For any given $g\geq3$, choose such a genus $G\geq g$.
Extend an actual boundary mapping class by the identity
to the larger surface. The corresponding Johnson
derivation extends by zero on the additional handles,
as in the actual geometric-image map of F.3.24.
For the inclusion $i:H_g\hookrightarrow H_G$, the
literal contraction formula gives
\[
                       R_G(i^{\otimes9}T)=iR_g(T).
                                                               \tag{RV.6}
\]
The left side is zero on the extended actual Johnson
image. Since $i$ is injective, (RV.6) gives (RV.1).
The strict Malcev-action identification used in F.3.24
transports the same assertion to its actual geometric
graded image. There is no assertion of density of a
discrete filtration kernel. Finally all maps and the
nonzero determinant are rational, so extension to
each characteristic-zero field $\Q_\ell$ preserves
the proof, including primes dividing the integer
determinant.
\end{proof}
